% s-320-p320-m002; original PDF 320, 321
\Aphorism{s-320-p320-m002}{377}
Hat man den ganzen Sinn und Gültigkeitsbezirk der „Wiedergeburt“ ermessen — \Sperr{Ich }bin zu irgend einer Zeit mit irgend einem Subjecte identisch —, so stellt sich noch einmal die ganze \Foreign{Neminem-laede}-Moral vor uns auf, als hedonistische Forderung des Egoismus selbst, eines auf „Alle“ erweiterten Egoismus. Sobald man aber diese Forderung schon für den gewöhnlichen Egoismus ablehnt und höhere Imperative als die blosse Schonung und Nichtschädigung zulässt, füllt mit jener Ausdehnung des Subjectbegriffs die letzte Schranke, und die menschlichen Dinge erscheinen in einer neuen Unschuld. Die grosse Freizügigkeit beginnt, und zu allen Zielen sind alle Wege offen; die Thatenmörder unter den Gedanken sind gerichtet, und das Leben wird zum unbegrenzten Experiment. Wir nannten Den heroisch, der sein empirisches Ich nicht schont; die Möglichkeit eines transcendenten Heroismus dämmert auf, der so philosophirt: An jeder meiner Thaten bin ich als Thäter, als Werkzeug, als Opfer betheiligt; indem ich sie thue, erkläre ich meine Bereitwilligkeit zu dieser Vielfachheit der Beziehungen, die alles Unrecht löscht. Mein ganzes Ich, mit allen seinen Möglichkeiten, soll leiden, wenn es nothwendig ist; ich handle nicht wesentlich anders, sondern nur beschränkter und wirkungsärmer, wenn ich mein \Sperr{augenblickliches Ich} zur ausschliesslichen Passivität verdamme und auf die Mitleidenschaft der \Sperr{augenblicklich Anderen} verzichte!
\par

